\documentclass[aspectratio=169, 10pt]{beamer}

\usetheme{Madrid}
\usecolortheme{default}
\definecolor{ucbblue}{RGB}{0, 51, 102}

\setbeamercolor{structure}{fg=ucbblue}
\setbeamercolor*{palette primary}{fg=white,bg=ucbblue}
\setbeamercolor*{palette secondary}{fg=white,bg=ucbblue!80}
\setbeamercolor*{palette tertiary}{fg=white,bg=ucbblue!90}
\setbeamercolor{title}{fg=white, bg=ucbblue}
\setbeamercolor{frametitle}{fg=white, bg=ucbblue}
\setbeamercolor{item}{fg=ucbblue}

\usepackage[T1]{fontenc}
\usepackage[utf8]{inputenc}
\usepackage{tgpagella}
\usefonttheme{serif}
\usepackage[spanish, es-noshorthands]{babel}
\usepackage{amsmath, amssymb, bm}
\usepackage{graphicx}
\usepackage{booktabs} % CRITICO PARA \toprule, \midrule, \bottomrule
\usepackage[most]{tcolorbox}
\usepackage{tikz}
\usepackage{pgfplots}
\pgfplotsset{compat=1.18}
\usepackage[american]{circuitikz}

\title[Transitorios RC y RL]{Unidad EC3: Respuesta Transitoria en el Dominio del Tiempo}
\subtitle{De Redes Estáticas a Circuitos Dinámicos de Primer Orden (RC/RL)}
\author[B. Quiroga Turdera]{M.Sc. Ing. Bernardo Quiroga Turdera}
\institute[UCB Tarija]{Universidad Católica Boliviana ``San Pablo'' — Sede Tarija}
\date{Semana 10}

\begin{document}

\begin{frame}
    \titlepage
\end{frame}

\begin{frame}{De Circuito Estático a Circuito Dinámico}
    \textbf{¿Qué nos entregó la Unidad EC2?}
    \begin{itemize}
        \item Análisis puramente resistivo en Estado Estacionario.
        \item Relaciones \textbf{algebraicas} ($V = I \cdot R$). Si una fuente cambia, el circuito responde de forma instantánea.
    \end{itemize}
    \vspace{0.3cm}
    \textbf{El Salto hacia EC3}
    \begin{itemize}
        \item Agregamos elementos que \textbf{almacenan energía}.
        \item Ya no basta con conocer el presente del circuito; importa su \textbf{condición inicial} (estado previo).
        \item Aparece una nueva variable explícita: \textbf{El Tiempo ($t$)}.
    \end{itemize}
    \vspace{0.3cm}
    \begin{center}
        \textbf{Línea Temporal de Conmutación:}\\
        $0^-$ (Antes) \quad $\rightarrow$ \quad $0^+$ (Inmediatamente después) \quad $\rightarrow$ \quad Transitorio \quad $\rightarrow$ \quad $\infty$ (Estado Final)
    \end{center}
\end{frame}

\begin{frame}{Elementos de Almacenamiento y Leyes Constitutivas}
    \begin{columns}
        \column{0.5\textwidth}
        \begin{center}
            \textbf{El Capacitor ($C$)}\\
            Almacena energía en un campo eléctrico.\\
            $w_C = \frac{1}{2} C v_C^2$
            \vspace{0.2cm}
            
            \begin{circuitikz}[scale=0.8, transform shape]
                \draw (0,0) to[C, l={$C$}, v={$v_C$}, i={$i_C$}] (2,0);
            \end{circuitikz}
            \vspace{0.2cm}
            
            $$i_C(t) = C \frac{dv_C(t)}{dt}$$
            \small \textit{La corriente depende de qué tan rápido cambie el voltaje en el tiempo.}
        \end{center}

        \column{0.5\textwidth}
        \begin{center}
            \textbf{El Inductor ($L$)}\\
            Almacena energía en un campo magnético.\\
            $w_L = \frac{1}{2} L i_L^2$
            \vspace{0.2cm}
            
            \begin{circuitikz}[scale=0.8, transform shape]
                \draw (0,0) to[L, l={$L$}, v={$v_L$}, i={$i_L$}] (2,0);
            \end{circuitikz}
            \vspace{0.2cm}
            
            $$v_L(t) = L \frac{di_L(t)}{dt}$$
            \small \textit{El voltaje depende de qué tan rápido cambie la corriente en el tiempo.}
        \end{center}
    \end{columns}
\end{frame}

\begin{frame}{Continuidad: La Inercia Eléctrica}
    Bajo condiciones ordinarias de circuitos (fuentes de energía finitas), la energía no se puede "teletransportar" ni cambiar en cero segundos. Esto impone reglas fundamentales:
    \vspace{0.3cm}
    
    \begin{block}{Regla de Continuidad del Capacitor (Inercia de Tensión)}
        El voltaje a través de un capacitor ideal no puede cambiar instantáneamente.
        $$v_C(0^+) = v_C(0^-)$$
    \end{block}
    
    \begin{block}{Regla de Continuidad del Inductor (Inercia de Corriente)}
        La corriente a través de un inductor ideal no puede cambiar instantáneamente.
        $$i_L(0^+) = i_L(0^-)$$
    \end{block}
    
    \vspace{0.2cm}
    \small \textbf{Nota de Cuidado:} No asuma que en $t=0^+$ el capacitor "siempre vale 0V". Eso \textbf{solo} es cierto si antes de conmutar ($0^-$) estaba completamente descargado.
\end{frame}

\begin{frame}{Circuito RC Serie: Del Esquema a la Ecuación Diferencial}
    Supongamos un capacitor \textbf{inicialmente descargado} ($v_C(0^-) = 0\,\text{V}$). En $t=0$, cerramos el interruptor.
    
    \begin{columns}[T] % Alineación superior para evitar desbordes
        \column{0.45\textwidth}
        \begin{center}
            % Escala reducida a 0.75 para evitar Overfull \hbox
            \begin{circuitikz}[scale=0.75, transform shape]
                \draw (0,0) to[V, l={$V_s=10\,\text{V}$}] (0,3) 
                      to[closing switch, l={$t=0$}] (2,3) 
                      to[R, l={$R=2\,\text{k}\Omega$}, v={$v_R$}] (4,3) 
                      to[C, l={$C=100\,\mu\text{F}$}, v={$v_C$}, i={$i_C$}] (4,0) -- (0,0);
            \end{circuitikz}
        \end{center}
        
        \column{0.55\textwidth}
        \small % Fuente reducida para evitar Overfull \vbox
        Aplicamos Ley de Voltajes de Kirchhoff (KVL):
        $$V_s - v_R - v_C = 0$$
        
        Sustituimos Ley de Ohm ($v_R = R \cdot i_C$):
        $$V_s - R \cdot i_C - v_C = 0$$
        
        Ley Constitutiva del Capacitor ($i_C = C \frac{dv_C}{dt}$):
        $$V_s - R \left( C \frac{dv_C}{dt} \right) - v_C = 0$$
        
        Reordenando, obtenemos la \textbf{EDO}:
        \begin{alertblock}{EDO de Primer Orden}
            \centering
            $RC \frac{dv_C}{dt} + v_C = V_s$
        \end{alertblock}
    \end{columns}
\end{frame}

\begin{frame}{Condiciones de Frontera y la Constante de Tiempo ($\tau$)}
    Para resolver la EDO: $0.2 \frac{dv_C}{dt} + v_C = 10$, necesitamos tres ingredientes clave:
    \vspace{0.2cm}
    
    \begin{enumerate}
        \item \textbf{Condición Inicial:} Por continuidad, sabemos que: \\
        $v_C(0^+) = v_C(0^-) = \mathbf{0\,\text{V}}$.
        
        \item \textbf{Condición Final (Estado Estacionario, $t \to \infty$):} \\ 
        Mucho tiempo después, el capacitor se carga y su voltaje deja de cambiar ($\frac{dv_C}{dt} \to 0$). Evaluando la EDO en $\infty$: \\
        $RC (0) + v_C(\infty) = 10 \implies v_C(\infty) = \mathbf{10\,\text{V}}$.
        
        \item \textbf{Constante de Tiempo ($\tau$):} \\
        Determina la escala de rapidez con la que reacciona el sistema: $\tau = R_{eq} C$. \\
        $\tau = (2\,\text{k}\Omega)(100\,\mu\text{F}) = \mathbf{0.2\,\text{s}}$.
    \end{enumerate}
    
    \vspace{0.2cm}
    \textit{Nota: $R_{eq}$ es la resistencia de la red pasiva (fuentes apagadas) vista desde el capacitor.}
\end{frame}

\begin{frame}{Forma General de la Respuesta Exponencial}
    Una vez obtenidos esos tres valores, no necesitamos resolver la EDO desde cero cada vez. Todo circuito de primer orden (con fuentes DC) tiene como solución:
    
    \begin{block}{Ecuación Universal de Primer Orden}
        $$x(t) = x(\infty) + \left[ x(0^+) - x(\infty) \right] e^{-t/\tau}$$
    \end{block}
    
    Lectura conceptual:
    $$\text{Respuesta}(t) = \text{Valor Final} + (\text{Desviación Inicial}) \cdot \text{Decaimiento Exponencial}$$
    
    Para nuestro ejemplo RC particular:
    $$v_C(t) = 10 + [0 - 10] e^{-t/0.2}$$
    $$v_C(t) = 10 (1 - e^{-5t}) \, \text{V}$$
\end{frame}

\begin{frame}{Lectura de la Curva y la Regla Práctica de los $5\tau$}
    \begin{columns}
        \column{0.5\textwidth}
        \begin{tikzpicture}
            \begin{axis}[
                width=\textwidth, height=5.5cm,
                xlabel={$t$ (s)}, ylabel={$v_C(t)$ (V)},
                xmin=0, xmax=1.2, ymin=0, ymax=11,
                grid=major,
                domain=0:1.2, samples=100,
                xtick={0, 0.2, 0.4, 0.6, 0.8, 1.0, 1.2},
                xticklabels={0, $1\tau$, $2\tau$, $3\tau$, $4\tau$, $5\tau$, $6\tau$}
            ]
            \addplot[blue, thick] {10*(1-exp(-x/0.2))};
            \addplot[dashed, red, thick] coordinates {(0,10) (1.2,10)};
            \addplot[mark=*, black] coordinates {(0.2, 6.32)} node[below right] {63.2\%};
            \addplot[mark=*, black] coordinates {(1.0, 9.93)} node[below right] {99.3\%};
            \end{axis}
        \end{tikzpicture}
        
        \column{0.5\textwidth}
        \textbf{Hitos de la Constante de Tiempo:}
        \begin{itemize}
            \item En $t = 1\tau$: La variable ha completado el $\approx 63.2\%$ de su transición total. ($e^{-1} \approx 0.368$ restante).
            \item En $t = 5\tau$: El error respecto al valor final es menor al $1\%$ ($e^{-5} \approx 0.0067$).
        \end{itemize}
        
        \vspace{0.3cm}
        \textbf{Convención de Tiempo de Asentamiento ($t_s$):}
        En ingeniería práctica, asumimos que el transitorio termina en:
        $$t_s \approx 5\tau$$
        (Para nuestro ejemplo: $1.0\,\text{s}$).
    \end{columns}
\end{frame}

\begin{frame}{Transferencia Conceptual: Del Circuito RC al RL}
    La misma arquitectura matemática rige a los inductores, pero la variable inercial es la corriente.
    \vspace{0.3cm}
    
    \begin{table}
        \centering
        \renewcommand{\arraystretch}{1.4}
        \resizebox{\textwidth}{!}{% Previene desbordamiento horizontal
        \begin{tabular}{lcc}
        \toprule
        \textbf{Característica} & \textbf{Circuito RC} & \textbf{Circuito RL} \\
        \midrule
        Variable de Estado Continua & $v_C(t)$ & $i_L(t)$ \\
        Relación Constitutiva & $i_C = C \frac{dv_C}{dt}$ & $v_L = L \frac{di_L}{dt}$ \\
        Condición de Continuidad & $v_C(0^+) = v_C(0^-)$ & $i_L(0^+) = i_L(0^-)$ \\
        Constante de Tiempo ($\tau$) & $R_{eq} \cdot C$ & $L / R_{eq}$ \\
        Estado Estacionario DC ($t \to \infty$) & $i_C \to 0$ (Circuito Abierto) & $v_L \to 0$ (Cortocircuito) \\
        \bottomrule
        \end{tabular}%
        }
    \end{table}
    \textit{Si aumentamos la resistencia en un RC, se vuelve más \textbf{lento}. Si aumentamos la resistencia en un RL, se vuelve más \textbf{rápido} ($\tau = L/R$).}
\end{frame}

\begin{frame}{Validación Próxima: Hacia la Herramienta Computacional}
    En la próxima fase, conectaremos el modelo analítico (matemático) con el modelo computacional (SPICE) y la medición física real.
    \vspace{0.4cm}
    
    \begin{center}
        \textbf{Modelo Analítico (EDO)} $\xrightarrow{\text{Compara}}$ \textbf{Simulación (}\texttt{.tran}\textbf{)} $\xrightarrow{\text{Compara}}$ \textbf{Osciloscopio}
    \end{center}
    
    \vspace{0.4cm}
    \textbf{Contrastes Esperados:}
    \begin{enumerate}
        \item ¿La curva física arranca en la condición inicial correcta de voltaje/corriente?
        \item ¿Se estabiliza en la asíntota ($t \to \infty$) calculada?
        \item ¿Tarda efectivamente $\sim 5\tau$ en asentarse en el osciloscopio?
    \end{enumerate}
\end{frame}

\end{document}
